\documentclass[11pt]{article}

\usepackage[margin=1in]{geometry}
\usepackage{hyperref}
\usepackage{booktabs}
\usepackage{amsmath}
\usepackage{amssymb}
\usepackage{graphicx}
\usepackage[numbers,sort&compress]{natbib}
\usepackage{xcolor}
\usepackage{array}

\hypersetup{
  colorlinks=true,
  linkcolor=blue!60!black,
  citecolor=blue!60!black,
  urlcolor=blue!60!black
}

\title{\textbf{Less Is More: Cognitive Load and the Single-Prompt\\
Ceiling in LLM Mathematical Reasoning}}

\author{
  Manuel Israel C\'{a}zares \\
  Bytepro AI \\
  Mazatl\'{a}n, Sinaloa, Mexico \\
  \texttt{hello@bytepro.ai} $|$ \texttt{israel.cazares@gmail.com}
}

\date{April 2026}

\begin{document}

\maketitle

% ─────────────────────────────────────────────
\begin{abstract}
We present a systematic empirical study of prompt engineering for
formal mathematical reasoning in the context of the SAIR Equational
Theories Stage 1 competition. The task requires deciding whether one
equational law implies another over all magmas (algebraic structures
with a single binary operation)~---~a problem that is undecidable in
general, where finite counterexample search is practically useful
for \textsc{False} despite being theoretically incomplete. Over
five weeks, we designed, tested, and analyzed more than~40 prompt
variants, ranging from~0 to~4{,}878 bytes, across four evaluation
splits and three language models (gpt-oss-120b, Llama~3.3~70B,
Gemma~4~31B).

Our central finding is an \textbf{empirical saturation region} within
the family of heuristic-guidance prompts we explored. Our two full-scale
($n=400$) evaluations anchor the central result: AN38 reaches 71.8\%
and AN45c reaches 79.3\% balanced hard accuracy for gpt-oss-120b,
compared to a~59.75\% no-cheatsheet baseline (95\%~CI: [54.9\%,
64.4\%]). Smaller-sample local evaluations shown in
Table~\ref{tab:single} span 54--74\%, with correspondingly
wider confidence intervals. Two considerations bear on this range
within our explored family. First, a theoretical limitation applies
to any prompt relying on finite counterexample enumeration: some
\textsc{False} instances may require counterexamples that are not
finite~\citep{tao2025etp}, though we have not established that this
explains the specific errors observed on hard3. Second, full-leaderboard
evidence associates rule-system content, not merely prompt length, with
degraded \textsc{True} recall on Llama~3.3~70B
(Section~\ref{sec:instability-llama}). Across our own prompt variants,
accuracy changes are fragile and non-monotonic; variants that differ in
ordering also differ in other respects, so we do not isolate an ordering
effect.

We also document a \textbf{distribution mismatch failure mode}: a
rule that appeared correct when validated on a \textsc{False}-heavy
subset (hard1, 35\% \textsc{True}) catastrophically failed on a balanced
subset (hard2, 50\% \textsc{True}), incorrectly blocking~51\% of
\textsc{True} problems. Our best submission (AN45c, 2{,}252~bytes) achieves \textbf{79.25\%}
accuracy on hard3 ($n=400$; 95\%~CI: [75.0\%, 82.9\%]), with
\textsc{True} recall of~95.9\% and \textsc{False} recall of~63.4\%,
representing a~+19.5 percentage-point improvement over the
no-cheatsheet baseline~(59.75\%; 95\%~CI: [54.9\%, 64.4\%]). A cross-provider validation run on
OpenRouter/DeepInfra~bf16 ($n=20$) yielded~90--95\%, consistent with
the full-scale result. The key design decision~---~placing the
trivial-magma check before the counterexample table~---~accounts for
the primary performance gain over its predecessor AN38~(71.8\%),
not the addition of new content. We release all prompt variants,
evaluation scripts, and results tables.

Post-submission validation against the SAIR official benchmark reveals
that gains within the saturation region are distribution-fragile: our
most engineered variant (AN45c) falls below the no-cheatsheet baseline
on the official evaluation while the simpler predecessor (AN38)
produces a robust improvement
(Section~\ref{sec:local-official}). Separately, full-leaderboard
evidence shows a directional research$\to$scored association under a
homogeneous-linear null (Section~\ref{sec:official}).

\textit{Note: This version incorporates the complete final competition leaderboard (310 evaluated submissions of 1{,}007 registered participants).}
\end{abstract}

% ─────────────────────────────────────────────
\section{Introduction}

Large language models have demonstrated surprising competence on
mathematical reasoning tasks, yet their behavior on problems requiring
formal logical completeness~---~where a single counterexample suffices
to disprove a claim~---~remains poorly understood. The SAIR Equational
Theories Stage~1 competition~\citep{tao2026sair} provides an unusually
clean testbed for studying this question: given two equations over magmas
(sets with a single binary operation), decide whether the first implies
the second universally. The problem is computationally asymmetric:
\textsc{False} instances are often certifiable by exhibiting a
small finite counterexample, while \textsc{True} instances require a
universal proof that no counterexample exists, including infinite
algebraic structures.

Unlike prompt engineering for standard benchmarks such as
GSM8K~\citep{cobbe2021gsm8k}, algebraic implication over magmas
exhibits an inherent asymmetry. While many positive implications can
be established through formal proof, negative implications may
require counterexample magmas that are not finite, making exhaustive
finite search insufficient in general~\citep{tao2025etp}~---~a
fundamentally different reasoning regime from standard benchmarks.

This asymmetry creates a natural design space for prompt engineering.
One might expect that providing a language model with a library of known
counterexamples would systematically improve \textsc{False} accuracy,
while a brief instruction about singleton-forcing equations would
improve \textsc{True} accuracy. Our experiments reveal a more complex
picture.

Over five weeks of systematic experimentation prior to the April~20,
2026 competition deadline, we tested more than~40 prompt variants on
over~1{,}000 labeled problems across four dataset splits. We found that
prompt complexity and multi-model generalization are inversely related:
the prompts that most improved gpt-oss-120b's performance on
\textsc{False}-heavy subsets (e.g., AN3c's 4{,}306-byte \textsc{Block}
system achieving~78.3\% on hard1) performed at or below baseline on
balanced subsets and collapsed to near-zero \textsc{True} recall on
Llama~3.3~70B. Conversely, the most compact effective prompt (AN19c,
289~bytes) performed within~2 percentage points of complex variants on
gpt-oss-120b while being the only variant that maintained meaningful
\textsc{True} recall on Llama.

The phenomenon we document within this prompt family~---~which we call
the \textbf{single-prompt ceiling}~---~is the upper-performing empirical
range achieved by our explored heuristic-guidance prompts on
gpt-oss-120b for this task. Within this family, our full-scale results
reached approximately 71--79\% balanced hard accuracy (AN38, AN45c),
with smaller-sample variants spanning a wider range down to~54\%
(Table~\ref{tab:single}); gpt-oss-120b achieves only~26.5\%
\textsc{False} recall with default reasoning and no cheatsheet on the
official benchmark (our own controlled baseline measurement
yields~38.0\% \textsc{False} recall under Together~AI inference; see
Section~\ref{sec:results}). We do not isolate which specific prompt
properties drive the gap between our upper- and lower-performing
variants within this family; Sections~\ref{sec:analysis}
and~\ref{sec:ceiling} discuss candidate factors and their limitations.

We make the following contributions:
\begin{enumerate}
  \item \textbf{Systematic ablation study}: 40+ prompt variants tested
    on labeled splits, enabling controlled analysis of which design
    choices matter.
  \item \textbf{Distribution mismatch failure mode}: Quantified how
    validation on \textsc{False}-heavy data produces incorrect
    conclusions when the target distribution is balanced.
  \item \textbf{Multi-model generalization analysis}: First systematic
    study of prompt portability across gpt-oss-120b, Llama~3.3~70B,
    and Gemma~4~31B on this task.
  \item \textbf{Prompt-variant comparison}: AN45c, which places the
    trivial-magma check before the counterexample table, outperforms
    AN38, which places the counterexample table first, by~$+7.5$
    percentage points at full scale ($n=400$; marginal 95\% Wilson CIs
    [75.0\%, 82.9\%] vs.\ [67.1\%, 75.9\%]). Because both variants are
    evaluated on the same 400 problems, we additionally report a
    continuity-corrected McNemar test on paired per-problem predictions
    ($\chi^2=9.14$, $p=0.0025$; $b=61$, $c=31$), providing evidence
    against symmetric disagreement. AN45c and AN38 differ in ordering
    and in other respects (Section~\ref{sec:method}); we hypothesize but
    do not isolate an ordering effect (Section~\ref{sec:analysis}).
    See Section~\ref{sec:local-official} for official-benchmark divergence.
  \item \textbf{Practical guidelines}: Minimal effective prompts for
    multi-model deployment, with analysis of why simpler beats complex.
  \item \textbf{Cross-distribution validation.} Official benchmark and
    leaderboard evidence show that local gains within the saturation
    region need not transfer across problem distributions
    (Sections~\ref{sec:local-official} and~\ref{sec:official}).
\end{enumerate}

The rest of the paper is organized as follows. Section~\ref{sec:background}
provides background on magmas, equational implication, and the
Equational Theories Project. Section~\ref{sec:related} reviews related
work. Section~\ref{sec:method} describes our methodology.
Sections~\ref{sec:results}--\ref{sec:multimodel} present results,
analysis, and multi-model generalization findings.
Section~\ref{sec:ceiling} characterizes the single-prompt ceiling
theoretically. Section~\ref{sec:official} presents post-submission
validation against the official benchmark and cross-distribution
analysis of the complete competition leaderboard.
Section~\ref{sec:instability} analyzes run-to-run verdict instability and its relation to prompt design.
Section~\ref{sec:conclusion} concludes.

% ─────────────────────────────────────────────
\section{Background}
\label{sec:background}

\subsection{Magmas and Equational Laws}

A \emph{magma} is a set~$M$ equipped with a single binary operation
${\star} : M \times M \to M$, closed under that operation and subject
to no further axioms. No associativity, commutativity, identity
element, or invertibility is assumed. An \emph{equational law} over a
magma is a universally quantified identity of the form
$t_1(x, y, \ldots) = t_2(x, y, \ldots)$, where $t_1$ and~$t_2$ are
terms built from variables and~$\star$. The law holds in a magma
$(M,\star)$ if the identity is satisfied for every assignment of
variables to elements of~$M$.

\subsection{Equational Implication}

Given two equational laws $E_1$ and~$E_2$, we say $E_1$ \emph{implies}
$E_2$ (written $E_1 \Rightarrow E_2$) if every magma that
satisfies~$E_1$ also satisfies~$E_2$. Deciding this relation is
computationally asymmetric. A \textsc{False} instance~---~where
$E_1 \not\Rightarrow E_2$~---~can be certified by exhibiting a single
finite counterexample magma in which~$E_1$ holds but~$E_2$ fails; finite
counterexample search is often practically useful for \textsc{False}
despite being theoretically incomplete. A \textsc{True}
instance requires establishing that no counterexample exists among all
magmas, including infinite ones, for which no general algorithm is
known. The implication problem over magmas is undecidable in general. Finite
model search is not complete even for \textsc{False}: there exist
implications that hold in all finite magmas but fail over infinite
ones, so a fruitless finite search certifies nothing.

\subsection{The Equational Theories Project}

The Equational Theories Project~\citep{tao2025etp} is a large-scale
collaborative effort to map the implication structure of equational laws
over magmas, formalized in Lean~4. The project has verified
approximately~4{,}694 distinct equational laws and established the
implication status of roughly~22 million equation pairs, producing the
largest formally verified database of algebraic implications to date.
This dataset provides both the mathematical foundation and the training
signal for the SAIR competition benchmark.

\subsection{Why the Task Is Hard for Language Models}

The asymmetry between \textsc{True} and \textsc{False} creates a
fundamental challenge for any fixed reasoning strategy. Producing a
valid \textsc{False} verdict requires constructing or recalling a
specific finite structure~---~a task amenable to lookup but not to
general inference. Producing a valid \textsc{True} verdict requires
establishing universal quantification over an infinite class of
structures~---~a task that exceeds any finite enumeration and for which
no sound and complete proof procedure exists in the general case. A
language model operating without external symbolic tools must
approximate both tasks within a single generation, using pattern
recognition over its training distribution as a substitute for proof
search. This approximation is the object of study in the present work.

% ─────────────────────────────────────────────
\section{Related Work}
\label{sec:related}

\paragraph{LLM mathematical reasoning benchmarks.}
GSM8K~\citep{cobbe2021gsm8k} and MATH~\citep{hendrycks2021math}
established word-problem and competition-mathematics benchmarks that
remain standard, but both assess reasoning over numeric domains with
human-interpretable intermediate steps. Tasks requiring universal
logical closure~---~deciding that a statement holds over \emph{all}
instances of a structure~---~are fundamentally different: a single
counterexample refutes a universal claim, yet no finite check can
confirm one. The SAIR Equational Theories benchmark evaluates~25 models
across~200 problems under varying conditions. The official results show
that even the strongest available model (Gemini~1.5~Pro) achieves~90.2\%
on hard problems without a cheatsheet, while weaker models cluster near
chance, making the benchmark unusually diagnostic of ceiling effects.

\paragraph{Prompt engineering for formal reasoning.}
Chain-of-thought prompting~\citep{wei2022cot} demonstrated that
instructing models to produce intermediate reasoning steps substantially
improves performance on multi-step problems. Subsequent work established
that few-shot exemplars~\citep{brown2020gpt3} and zero-shot
chain-of-thought instructions~\citep{kojima2022zeroshot} are broadly
effective. However, these findings are primarily established on numeric
and commonsense domains. Less is known about how prompt complexity
interacts with formal algebraic reasoning, where rule fidelity~---~not
just step count~---~determines correctness.

\paragraph{Prompt complexity and instruction-following reliability.}
\citet{sharma2023sycophancy} documented that RLHF-trained models
consistently exhibit sycophancy across varied free-form generation
tasks, adjusting their stated conclusions toward what they perceive the
user wants even when incorrect. Within our prompt family, we observe a
related but distinct pattern specifically on Llama~3.3~70B: under
certain rule-system content, it can default to a near-constant
\textsc{False} output (Section~\ref{sec:multimodel}), a pattern better
predicted by rule-system content than by prompt length alone
(Section~\ref{sec:instability-llama}). Whether this reflects the same
underlying mechanism as sycophancy is untested; we note it here as a
related but distinct phenomenon. \citet{shi2023distracted} similarly
showed that irrelevant context in math problems degrades accuracy,
consistent with the general observation that additional prompt content
does not uniformly help reasoning, though the mechanism may differ
from ours.

\paragraph{The Equational Theories Project.}
The mathematical foundations of this competition derive from ongoing
work on automated proof search over equational theories~\citep{tao2025etp}.
The project established that implications between equational laws over
magmas exhibit complex dependency structures not tractable by
brute-force enumeration, motivating the use of language models as
approximate reasoners over this space.

\paragraph{In-context learning: length versus performance.}
\citet{liu2023lost} showed that model attention is non-uniform over
context: information at the beginning and end of a prompt is retrieved
more reliably than information in the middle. This has direct
implications for structured cheatsheets. Our finding that
trivial-magma-first ordering (AN45c) outperforms
counterexample-table-first ordering (AN38) by~$+7.5$ percentage points
at full scale ($n=400$; Section~\ref{sec:results}). We hypothesize that
placing the
trivial-magma check first primes the model's attention toward
\textsc{True} verdicts before engaging the CE search, with the first
substantive rule receiving disproportionate weight during generation;
this mechanism remains to be verified through attention analysis.

% ─────────────────────────────────────────────
\section{Methodology}
\label{sec:method}

\subsection{Task and Benchmark}

The SAIR Equational Theories Stage~1 competition requires deciding,
for a given pair of equations $(E_1, E_2)$ over magmas, whether
$E_1 \Rightarrow E_2$ holds universally (label: \textsc{True}) or
whether a counterexample magma exists (label: \textsc{False}). The
official evaluation judge is described in~\citet{sair2026judge}. All
experiments use the publicly available dataset\\
\texttt{SAIRfoundation/equational-theories-selected-problems}
(HuggingFace). We work with four labeled splits, summarized in
Table~\ref{tab:splits}.

\begin{table}[h]
\centering
\caption{Dataset splits used in this study.}
\label{tab:splits}
\begin{tabular}{lrrrl}
\toprule
Split & $n$ & \textsc{True} & \textsc{False} & Notes \\
\midrule
normal & 1{,}000 & 500 (50\%) & 500 (50\%) & Balanced; mostly \texttt{x = RHS} form \\
hard1  & 69      & 24 (35\%)  & 45 (65\%)  & \textsc{False}-heavy; dense nesting \\
hard2  & 200     & 100 (50\%) & 100 (50\%) & Balanced; structurally complex \\
hard3  & 400     & 195 (49\%) & 205 (51\%) & Near-balanced; primary eval split \\
\bottomrule
\end{tabular}
\end{table}

Hard3 serves as the primary evaluation split: its near-balanced
distribution and size ($n=400$) most closely approximate the
competition's private evaluation set. Hard1 is used selectively to
test \textsc{False}-detection strategies; hard2 provides a secondary
balanced check. Normal problems confirm that interventions do not
regress standard performance.

\subsection{Models and Inference Configuration}

We evaluate across three models matching the competition's official
multi-model setup:

\textbf{gpt-oss-120b} (primary): An open-weight 117B-parameter
Mixture-of-Experts model released by OpenAI under Apache~2.0
license~\citep{openai2025gptoss}, accessed via Together~AI~bf16 for local full-scale evaluations
($n=400$) and via OpenRouter/DeepInfra~bf16 for cross-provider
validation, the latter matching the official competition
pipeline~\citep{sair2026judge}
(\texttt{openai/gpt-oss-120b}). This model uses an
extended reasoning mode that produces chain-of-thought prior to its
final verdict. All competition-credit evaluations use this model.

\textbf{Llama~3.3~70B}: Accessed via Together~AI
(\texttt{meta-llama/Llama-3.3-70B-Instruct-Turbo}). Standard
instruction-tuned mode; no extended reasoning.

\textbf{Gemma~4~31B}: Accessed via Together~AI
(\texttt{google/gemma-4-31b-it}) with \texttt{max\_tokens=8192} to
enable the model's native reasoning trace. OpenRouter routing for this
model suppresses reasoning mode, causing it to default to a
near-constant \textsc{True} output ($\approx$53\%); all reported Gemma
results use Together~AI exclusively.

All inference runs use \texttt{temperature=0} and \texttt{seed=42}.
\texttt{max\_tokens} is set to~4{,}096 for gpt-oss-120b and~8{,}192
for Gemma~4~31B. A preliminary experiment established that truncation
at lower token budgets (512, 1{,}024, 2{,}048) produced 100\%
truncation-caused errors; 4{,}096 was the minimum budget at which
genuine reasoning errors first appeared. Estimated cost per problem is
\$0.005--\$0.01 depending on prompt length and model.

\subsection{Evaluation Metrics}

We report three primary metrics:
\textbf{Accuracy} (fraction correct),
\textbf{\textsc{True} recall} (fraction of \textsc{True}-labeled
problems answered \textsc{True}), and
\textbf{\textsc{False} recall} (fraction of \textsc{False}-labeled
problems answered \textsc{False}).
For multi-model comparisons we report the \textbf{3-model average}:
the unweighted mean of each model's accuracy on the same split, aligned
with the competition's scoring rule.

Non-determinism is a practical concern at \texttt{temperature=0}: we
observed up to $\pm$3 percentage points of variance across identical
runs (e.g., AN3d: 73.9\% vs.\ 79.7\% on separate executions). We
report observed values without averaging across runs, and flag cases
where $n \leq 20$ makes estimates unreliable ($\pm$10pp at the 95\%
level).

\subsection{Prompt Design Space}

Each submission is a single UTF-8 text file of at most~10KB containing
two placeholders, \texttt{\{\{ equation1 \}\}} and
\texttt{\{\{ equation2 \}\}}. We designed and evaluated~45+ prompt
variants over five weeks (AN-series: AN1 through AN45d).

\smallskip
\noindent\textbf{Pipeline note.} AN45c uses a self-contained template
format with \texttt{\{\{ equation1 \}\}} and \texttt{\{\{ equation2 \}\}}
placeholders inside the prompt body, requiring the \texttt{--raw-prompt}
flag for correct substitution. An early run omitted this flag; the
placeholders reached the model unsubstituted, producing an artifactual
56\% result that was identified and discarded. All other variants in
Table~\ref{tab:single} use \texttt{build\_prompt()}, which embeds
equations inline and is unaffected by this flag. The corrected AN45c
results (April~14, 2026) are the only ones reported here.
\smallskip

Variants differ along five design dimensions:

\begin{enumerate}
  \item \textbf{Counterexample (CE) table content}: which small magmas
    are provided, from~4 size-2 structures (early variants) to~7
    size-2 and~5 size-3 structures (AN38/AN45c).
  \item \textbf{Singleton-forcing rules}: heuristics for detecting
    equations that force all elements equal (TRUE shortcuts). S1 is
    sound; S2 is unsound for right-zero magmas (Section~\ref{sec:analysis}).
  \item \textbf{Block rules}: conditional classifiers for structural
    patterns predictive of \textsc{False}. The most aggressive variant
    (Block~1 in AN3c) blanket-classified self-referential patterns as
    \textsc{False}, achieving high recall on hard1 but catastrophic
    precision loss on balanced splits.
  \item \textbf{Instruction ordering}: whether the trivial-magma check
    or the CE table appears first. AN45c places the trivial-magma
    check first; AN38 places the CE table first.
  \item \textbf{Prompt length}: 0~bytes (baseline) to~4{,}878~bytes
    (AN5, the longest variant tested).
\end{enumerate}

Variants were selected for full-scale evaluation based on
small-sample performance ($n=20$--50), with the most promising
candidates validated at $n=200$--400. All prompt variants, evaluation
scripts, and results are available at our companion
repository~\citep{cazares2026sair}.

% ─────────────────────────────────────────────
\section{Results}
\label{sec:results}

We present results along four axes: single-model performance across
prompt variants (\S\ref{subsec:single}), cross-model generalization
(\S\ref{subsec:cross}), cross-dataset stability
(\S\ref{subsec:dataset}), and local vs.\ official benchmark divergence
(Section~\ref{sec:local-official}).

\paragraph{Sample size note.}
Results for AN45c on gpt-oss-120b are based on $n=400$ hard3
problems (full-scale corrected pipeline). Cross-model runs
(Llama~3.3~70B, Gemma~4~31B) and the DeepSeek~V3.2 result
use $n=50$, $n=20$, and $n=10$ respectively, as noted in
Table~\ref{tab:single}.

\subsection{Single-Model Results (gpt-oss-120b, hard3)}
\label{subsec:single}

Table~\ref{tab:single} reports accuracy, \textsc{True} recall,
\textsc{False} recall, and prompt size for all variants evaluated on
gpt-oss-120b on the hard3 split ($n=50$ unless noted), ordered by
descending accuracy. The no-cheatsheet baseline (own run, April~14,
2026, $n=400$) achieves~59.75\% overall (95\%~CI: [54.9\%, 64.4\%]),
with a severe \textsc{True} bias: 82.6\% \textsc{True} recall but
only~38.0\% \textsc{False} recall. This profile is consistent with
the structural \textsc{True} bias documented in the official benchmark
(89.2\% \textsc{True}, 26.5\% \textsc{False} on the hard split).

\begin{table}[h]
\centering
\caption{Single-model results on hard3 (gpt-oss-120b, $n=50$ unless noted).}
\label{tab:single}
\small
\begin{tabular}{lrrrr>{\raggedright\arraybackslash}p{4.2cm}}
\toprule
Variant & Acc & T\% & F\% & Bytes & Strategy \\
\midrule
AN45c ($n=400$)\textsuperscript{$\dagger$}
        & \textbf{79.3} & 95.9 & 63.4 & 2{,}252 & Trivial magma first + CE tables A--N \\
AN45c ($n=20$, OR)\textsuperscript{$\S$}
        & 90.0--95.0 & --- & --- & 2{,}252 & Cross-provider validation (OpenRouter/DeepInfra) \\
AN38 ($n=400$)\textsuperscript{$\ddagger$}
        & \textbf{71.8} & 78.5 & 65.4 & 1{,}776 & CE tables A--N (3-elem focus) \\
AN38 ($n=50$)  & 74.0 & 70.8 & 76.9 & 1{,}776 & CE tables A--N (3-elem focus) \\
AN43           & 72.0 & 54.2 & 88.5 & 2{,}171 & Controller arch (BLOCKs as router) \\
AN35           & 72.0 & 58.3 & 84.6 & 1{,}545 & 3-elem CE focus \\
AN35b          & 72.0 & 79.2 & 65.4 & 1{,}802 & Cautious TRUE + 3-elem CE \\
AN45d          & 70.0 & 100.0 & 40.0 & 2{,}538 & AN45c + corrected STEP~1 flag \\
AN36           & 70.0 & 50.0 & 88.5 & 1{,}205 & Aggressive FALSE prior \\
AN39           & 70.0 & 91.7 & 50.0 & 385   & Power-level prior \\
Baseline ($n=400$) & 59.75 & 82.6 & 38.0 & 0 & No cheatsheet (own run, April~14, 2026) \\
AN3c ($n=50$)  & 64.0 & 45.8 & 80.8 & 4{,}306 & BLOCK system \\
AN10           & 64.0 & ---  & ---  & 3{,}303 & Symbolic engine \\
AN19c          & 62.0 & 91.7 & 34.6 & 289   & Trivial magma hint only \\
AN42           & 62.0 & ---  & ---  & ---   & KB rewriting \\
AN40           & 60.0 & ---  & ---  & ---   & Semantic invariants \\
AN41           & 58.0 & ---  & ---  & ---   & Tamari/tree-structure \\
AN5            & 54.0 & ---  & ---  & 4{,}878 & Maximum-length CE table \\
\bottomrule
\multicolumn{6}{p{\linewidth}}{\small\textsuperscript{$\dagger$}Corrected pipeline (\texttt{--raw-prompt}); $n=400$; primary result.
\textsuperscript{$\ddagger$}Full-scale run; most reliable pre-fix estimate.
\textsuperscript{$\S$}Cross-provider check only; $n=20$, $\pm$10pp CI.} \\
\multicolumn{6}{p{\linewidth}}{\small\textit{Wilson 95\% confidence intervals:}
AN45c $n=400$, 317/400=79.3\%: \textbf{[75.0\%, 82.9\%]};
AN38 $n=400$, 287/400=71.8\%: [67.1\%, 75.9\%];
AN3c hard1 $n=69$, 54/69=78.3\%: [67.2\%, 86.4\%].}
\end{tabular}
\end{table}

Performance is non-monotonic in prompt length: the longest variant
(AN5, 4{,}878~bytes) is the worst-performing cheatsheet, while the
shortest effective variant (AN39, 385~bytes) ties with mid-length
prompts at~70\%. \textsc{True} and \textsc{False} recall trade off
sharply: no variant simultaneously achieves both above~80\% on hard3.
AN19c maximizes \textsc{True} recall~(91.7\%) at the cost of
near-complete \textsc{False} recall collapse~(34.6\%), while AN36
inverts this profile (50.0\% \textsc{True}, 88.5\% \textsc{False}).

The AN45d result (100\% \textsc{True}, 40\% \textsc{False}) warrants
attention: a minor modification to AN45c caused six
\textsc{False}$\to$\textsc{True} regressions, dropping overall
accuracy from~90\% to~70\%. The ordering effect in AN45c depends on
the exact token sequence, not on the semantic content of the rules.

\paragraph{Cheatsheet effect on model bias.}
The no-cheatsheet baseline exhibits strong structural \textsc{True}
bias: 82.6\% \textsc{True} recall versus only~38.0\% \textsc{False}
recall ($n=400$, own run). AN45c does not merely raise overall
accuracy~---~it rebalances this bias. \textsc{False} recall improves
by~$+25.4$ percentage points (38.0\%~$\to$~63.4\%), while
\textsc{True} recall increases further to~95.9\%. The cheatsheet
therefore serves two distinct functions: it corrects the model's
structural tendency to classify hard problems as \textsc{True}, and
it provides the finite counterexample evidence needed to produce
confident \textsc{False} verdicts.

\subsection{Cross-Model Results}
\label{subsec:cross}

Table~\ref{tab:multimodel} reports results for key variants across all
three competition models.

\begin{table}[h]
\centering
\caption{Multi-model results on hard3 ($n=20$ for AN45c; $n=50$ otherwise).}
\label{tab:multimodel}
\small
\begin{tabular}{lrrrrl}
\toprule
Variant & Bytes & gpt-oss & Llama & Gemma & 3-avg \\
\midrule
AN45c~-- Scenario~A (official)\textsuperscript{$\dagger$}
    & 2{,}252 & 95\% & 55\% & 53\% (OR) & \textbf{67.7\%} \\
AN45c~-- Scenario~B
    & 2{,}252 & 90\% & 55\% & 85\% (TAI) & 76.7\% \\
AN45c~-- Scenario~C
    & 2{,}252 & 95\% & 55\% & 85\% (TAI) & 78.3\% \\
AN38    & 1{,}776 & 74\% & 52\% & 54\% & 59.3\% \\
AN19c   & 289   & 62\% & \textbf{60\%} & 55\% & 59.0\% \\
Baseline & 0    & 59.75\% & 52\% & $\approx$50\% & --- \\
AN3c    & 4{,}306 & 64\% & $\approx$0\%\textsuperscript{$\ddagger$} & --- & --- \\
\bottomrule
\multicolumn{6}{p{0.92\textwidth}}{%
\textsuperscript{$\dagger$}Scenario~A uses OR/Novita~bf16 Gemma (reasoning suppressed,
near-constant \textsc{True} output). Conservative official lower bound.
Scenarios~B/C use Together~AI (8{,}192~tok, reasoning enabled); differ
only in GPT result (90\% vs.\ 95\%, $n=20$).
\textsuperscript{$\ddagger$}0\% \textsc{True} recall; overall accuracy reflects only \textsc{False} correct answers.}
\end{tabular}
\end{table}

We report~67.7\% (Scenario~A) as the conservative official 3-model
average. Scenarios~B and~C (76.7--78.3\%) represent Gemma performance
when correctly configured; we report them to distinguish configuration
artifact from model capability.

The most striking cross-model finding is the AN3c collapse on Llama:
the 4{,}306-byte \textsc{Block} system causes Llama to output
\textsc{False} for every problem, yielding 0\% \textsc{True} recall.
AN19c (289~bytes) is the only variant that produces balanced recall on
Llama (37.5\% \textsc{True}, 80.8\% \textsc{False}), and the only
variant where Llama marginally outperforms gpt-oss-120b (60\%
vs.~62\%, within noise).

\subsection{Cross-Dataset Generalization}
\label{subsec:dataset}

Table~\ref{tab:crossdataset} shows performance across splits for AN3c
and AN38~---~variants optimized on hard1 and hard3 respectively.

\begin{table}[h]
\centering
\caption{Cross-dataset performance for selected variants (gpt-oss-120b).}
\label{tab:crossdataset}
\begin{tabular}{llrrrr}
\toprule
Variant & Split & $n$ & Acc & T\% & F\% \\
\midrule
AN3c & hard1  & 69  & \textbf{78.3} & 66.7 & 84.4 \\
AN3c & hard2  & 50  & 60.0 & 50.0 & 69.2 \\
AN3c & hard3  & 50  & 64.0 & 45.8 & 80.8 \\
AN3c & normal & 100 & 92.0 & 89.3 & 95.5 \\
\midrule
AN38 & hard1  & 69  & 50.7 & 75.0 & 37.8 \\
AN38 & hard2  & 200 & 51.5 & 78.0 & 25.0 \\
AN38 & hard3  & 400 & \textbf{71.8} & 78.5 & 65.4 \\
AN38 & normal & 200 & \textbf{92.0} & 89.9 & 94.5 \\
\bottomrule
\end{tabular}
\end{table}

AN3c's Block~1 rule was developed against hard1's \textsc{False}-heavy
distribution (35\% \textsc{True}). On hard1 it achieves~78.3\%, the
highest single-split result in the study. On hard2 (balanced, 50\%
\textsc{True}), the same rule drops to~60.0\%~---~only~8 percentage
points above chance~---~because Block~1 misclassifies~51\% of
\textsc{True} problems as \textsc{False}. This 18.3 percentage-point
degradation illustrates the distribution mismatch failure mode
precisely. AN38 exhibits the complementary pathology: well-calibrated
for hard3~(71.8\% full-scale) but near chance on hard1~(50.7\%). Both
variants achieve $\approx$92\% on normal problems, confirming that
normal-split results are uninformative for distinguishing hard-problem
strategies.

% ─────────────────────────────────────────────
\subsection{Local vs.\ Official Benchmark Divergence}
\label{sec:local-official}

Table~\ref{tab:official_hard3} compares local evaluation results
against the SAIR official benchmark on hard3 (the competition's
reference format, as confirmed by the official smoke-test file
\texttt{problems\_hard3\_20.jsonl}~\citep{sair2026judge}).

\begin{table}[h]
\centering
\caption{Local vs.\ official benchmark results on hard3 (GPT-OSS 120B).
Official benchmark: $n=20$, DeepInfra~bf16. Local: $n=400$,
Together~AI~bf16. Baseline: no cheatsheet.}
\label{tab:official_hard3}
\small
\begin{tabular}{lrrrr}
\toprule
Variant & Local Acc. & Official Acc. & $\Delta$ vs.\ baseline\textsuperscript{a} & Official F1 \\
\midrule
Baseline & 59.75\% & 56.3\% & --- & $\approx$53\% \\
AN38     & 71.8\%  & 65.3\% & \textbf{+5.6pp} & \textbf{67.6\%} \\
AN45c    & \textbf{79.25\%} & 55.5\% & $-$4.3pp & 63.7\% \\
\bottomrule
\multicolumn{5}{l}{\textsuperscript{a}Delta relative to official baseline (56.3\%).}
\end{tabular}
\end{table}

AN38 produces a robust $+$5.6pp improvement over the official baseline
on hard3, consistent with its local result (71.8\%). AN45c, despite
achieving 79.25\% locally ($n=400$), scores 55.5\% on the official
benchmark~---~4.3pp \emph{below} the no-cheatsheet baseline.

\begin{table}[h]
\centering
\caption{Official benchmark results on hard2 (GPT-OSS 120B, $n=20$).
Hard2 baseline: 56.3\% accuracy, F1 $\approx$64.9\%.}
\label{tab:official_hard2}
\small
\begin{tabular}{lrrr}
\toprule
Variant & Official Acc. & $\Delta$ vs.\ baseline & Official F1 \\
\midrule
AN38  & 41.0\% & $-$15.3pp & 47.8\% \\
AN45c & 48.0\% & $-$8.3pp  & \textbf{61.2\%} \\
\bottomrule
\end{tabular}
\end{table}

On hard2, both variants fall below the no-cheatsheet baseline. AN38
incurs a larger accuracy penalty ($-$15.3pp) and F1 degradation
($-$17.1pp), while AN45c is less damaging ($-$8.3pp accuracy,
$-$3.7pp F1). Neither cheatsheet improves over baseline on hard2,
confirming that both are calibrated toward the hard3 distribution.
On the official hard2--hard3 pair, AN38's cross-split accuracy gap is
24.3pp (41.0\% vs.\ 65.3\%) and AN45c's is 7.5pp (48.0\% vs.\ 55.5\%).

% ─────────────────────────────────────────────
\section{Analysis}
\label{sec:analysis}

\subsection{Why AN45c Works: The Trivial Magma as an Exit Gate}

AN45c's primary mechanism is structural rather than informational.
STEP~1 instructs the model to check whether~$E_1$ contains a variable
appearing exactly once on one side~---~a condition that forces all
elements of any satisfying magma to be equal, collapsing it to the
trivial one-element structure and making~$E_2$ vacuously true. When
triggered, STEP~1 bypasses STEP~2 (the counterexample search) entirely
and commits to \textsc{True} before the CE~tables introduce
\textsc{False}-directional pressure.

On hard3, AN45c achieves~95.9\% \textsc{True} recall compared to the
baseline's~82.6\%~---~a gain of~$+13.3$ percentage points~---~while
achieving~63.4\% \textsc{False} recall versus the baseline's~38.0\%,
a gain of~$+25.4$ percentage points.
The overall improvement is~$+19.5$~pp at $n=400$ (corrected pipeline).
This gain appears to arise from reordering existing components rather
than adding new content: AN38 contains the same STEP~1 logic but places
it \emph{after} the CE table, producing~78.5\% \textsc{True}
recall and~65.4\% \textsc{False} recall at full scale ($n=400$).
The ordering difference thus manifests primarily as a
\textsc{True}-recall gap ($95.9\%$ vs.\ $78.5\%$) at a small
\textsc{False}-recall cost ($63.4\%$ vs.\ $65.4\%$). We hypothesize
that placing the trivial-magma check first primes the model's attention
toward \textsc{True} verdicts before the CE search introduces
\textsc{False}-directional pressure~---~though isolating this ordering
effect from confounds requires controlled experiments beyond the scope
of this study.

The same STEP~1 gate is a candidate explanation for AN45c's
local--official inversion (Section~\ref{sec:local-official}): despite
79.25\% local accuracy ($n=400$), AN45c scores 55.5\% on official
hard3~---~4.3pp below the no-cheatsheet baseline. We hypothesize that
problems satisfying the singleton test are committed to \textsc{True}
before counterexample evidence is considered, generating forced errors
under distribution shift---compatible with the distribution-mismatch
failure mode documented for AN3c on hard1/hard2
(Section~\ref{sec:results}). This remains a hypothesis: isolating the
gate from other AN45c--AN38 differences would require controlled
ablations beyond the present study.

\subsection{Why the Merge Ceiling Holds}

The AN38 variant was constructed by merging two complementary
predecessors: AN35~(\textsc{True}=58.3\%, \textsc{False}=84.6\%) and
AN35b~(\textsc{True}=79.2\%, \textsc{False}=65.4\%). The result was
AN38: \textsc{True}=70.8\%, \textsc{False}=76.9\%~---~a near-arithmetic
mean of the two parents~---~not a combination of their strengths. Three
subsequent merge attempts (AN45d, AN45e, AN45f) all produced accuracy
at or below the mean. A static instruction set cannot simultaneously
prime two competing inference strategies: when the model encounters
rules favoring both labels, it averages rather than selects.

\subsection{The Token Cap Finding}

Gemma~4~31B exhibits a distinctive failure mode when \texttt{max\_tokens}
falls below the budget required to complete CE verification. At~2{,}048
tokens, Gemma with AN45c produces~50\% accuracy with~0\%
\textsc{False} recall~---~effectively outputting \textsc{True} for
every problem. At~8{,}192 tokens, the same model and prompt produces
85\% accuracy (\textsc{True}=90\%, \textsc{False}=80\%). The mechanism:
Gemma exhausts its token budget during STEP~2, fails to produce a
\textsc{Counterexample} block, and defaults to the STEP~1 exit gate
verdict (\textsc{True}). This is not a model capability failure; it is
a configuration artifact. Reported Gemma results are only valid when
the token budget is sufficient to complete the full reasoning trace.

\subsection{Why Theoretical Approaches Fail}

Four variants introducing explicit mathematical reasoning frameworks all
degraded performance relative to the 68\% baseline ($n=50$ subset,
consistent with Table~\ref{tab:theory} deltas):

\begin{table}[h]
\centering
\caption{Theory-based variants on hard3 (gpt-oss-120b, $n=50$).}
\label{tab:theory}
\begin{tabular}{llrr}
\toprule
Variant & Framework & Acc & $\Delta$ vs.\ baseline \\
\midrule
AN39 & Power-level prior       & 70\% & $+$2pp \\
AN42 & Knowledge-base rewriting & 62\% & $-$6pp \\
AN40 & Semantic invariants      & 60\% & $-$8pp \\
AN41 & Tamari/tree-structure    & 58\% & $-$10pp \\
\bottomrule
\end{tabular}
\end{table}

The consistent direction of failure is \textsc{False}$\to$\textsc{True}
error increase. When given a structured reasoning framework, the model
generates plausible arguments satisfying the framework's criteria and
commits to \textsc{True}, even for problems where a small
counterexample exists. False precision in the instructed framework is
worse than no framework at all.

\subsection{A Note on Numerical Coincidence}

AN45c's conservative 3-model average (Scenario~A: 67.7\%) and AN3c's
hard1 single-model accuracy~(78.3\%) are occasionally cited together in
competition discussions. They are not comparable: the former is a
balanced hard3 average across three models under a specific provider
configuration; the latter is a single-model result on a
\textsc{False}-heavy split where a blanket-\textsc{False} strategy
would already score~65\%. AN3c on a balanced split achieves~60--64\%.

\paragraph{Structural classification hypothesis.}
Our results suggest that large language models in this task behave not
as symbolic theorem provers but as heuristic classifiers over algebraic
structure patterns. The prompt encodes a decision boundary over
structural features~---~variable repetition, nesting depth, singleton
forcing~---~rather than a proof system. This framing explains both the
effectiveness of counterexample tables (which encode explicit structural
decision rules) and the failure of procedural approaches (which require
genuine symbolic execution). We term this the \emph{router hypothesis}:
the model routes problems to cached structural patterns rather than
deriving answers compositionally. This has a direct practical implication:
prompt improvements that add new structural rules can help, but only up
to the capacity of the model to maintain and apply those rules
consistently~---~which is precisely the saturation boundary we observe.

% ─────────────────────────────────────────────
\section{Multi-Model Generalization}
\label{sec:multimodel}

\subsection{AN19c as the Model-Agnostic Minimum}

AN19c (289~bytes) consists of three natural-language hints with no
counterexample tables: a reminder that one-element magmas trivially
satisfy all equations, a note that equations with a free variable
isolated on one side are likely \textsc{True}, and an instruction to
output a verdict with brief reasoning. On gpt-oss-120b it achieves~62\%
on hard3~---~below most CE-table variants~---~but it is the only prompt
that produces meaningful \textsc{True} recall on Llama~3.3~70B (37.5\%
\textsc{True}, 80.8\% \textsc{False}, 60\% overall) and competitive
performance on Gemma~4~31B (55\%, Together~AI).

Table~\ref{tab:multimodel4} summarizes multi-model performance for
variants with coverage across three or more models.

\begin{table}[h]
\centering
\caption{Multi-model generalization on hard3 ($n=10$--50; see notes).}
\label{tab:multimodel4}
\small
\begin{tabular}{lrrrrrr}
\toprule
Variant & Bytes & gpt-oss & Llama & Gemma (TAI) & DeepSeek & 4-avg \\
\midrule
AN19c   & 289   & 62\% & \textbf{60\%} & 55\% & \textbf{80\%}\textsuperscript{$\dagger$} & $\approx$64.3\% \\
AN38    & 1{,}776 & 74\% & 52\% & 54\% & 60\%\textsuperscript{$\dagger$} & $\approx$60.0\% \\
AN45c   & 2{,}252 & 90--95\% & 55\% & \textbf{85\%} & --- & --- \\
Baseline & 0   & 59.75\% & 52\% & $\approx$50\% & --- & --- \\
AN3c    & 4{,}306 & 64\% & $\approx$0\%\textsuperscript{$\ddagger$} & --- & --- & --- \\
\bottomrule
\multicolumn{7}{l}{\textsuperscript{$\dagger$}DeepSeek $n=10$; $\pm$15pp CI.
  \textsuperscript{$\ddagger$}0\% \textsc{True} recall.}
\end{tabular}
\end{table}

\subsection{Model-Specific Bias Profiles}

\textbf{gpt-oss-120b} is \textsc{True}-biased at baseline: 82.6\%
\textsc{True} recall versus~38.0\% \textsc{False} recall on hard3. CE
table prompts partially correct this but never eliminate the bias~---~even
AN38's best full-scale result~(71.8\%) reduces the baseline's 44.6-point
recall gap (82.6\% \textsc{True} vs.\ 38.0\% \textsc{False}) to just
13.1~points (78.5\% \textsc{True} vs.\ 65.4\% \textsc{False})~---~a
substantial improvement, but the imbalance persists.

\textbf{Llama~3.3~70B} exhibits the opposite profile: near-baseline
\textsc{False} recall~(80.8\% with AN19c) but severely depressed
\textsc{True} recall (37.5\% with the best prompt; effectively~0\%
with our CE-table prompt family). Within our design family, longer
prompts collapsed \textsc{True} recall further; however,
Section~\ref{sec:instability-llama} shows with full-leaderboard data
that the collapse is driven by prompt content rather than length.

\textbf{Gemma~4~31B} is token-gated rather than directionally biased:
its failure mode is token budget exhaustion, which triggers the STEP~1
exit gate and produces near-constant \textsc{True} output. With
sufficient budget~(8{,}192~tokens), Gemma's recall profile is balanced
and strong~(85\% with AN45c on $n=20$).

\subsection{DeepSeek V3.2: Cost Efficiency Signal}

DeepSeek~V3.2 with AN19c achieves~80\% on hard3 ($n=10$) at
$\approx$\$0.0008 per problem~---~roughly one order of magnitude cheaper
than gpt-oss-120b at comparable accuracy. With AN38~(1{,}776~bytes),
DeepSeek drops to~60\%, confirming the pattern observed in Llama: CE
table prompts hurt capable reasoners by introducing distracting
structure that interferes with internal algebraic reasoning. The AN19c
result on DeepSeek is the most cost-efficient signal in the study
($n=10$; highest-priority target for scale validation).

% ─────────────────────────────────────────────
\section{The Single-Prompt Ceiling}
\label{sec:ceiling}

\subsection{Formal Characterization}

We define the \emph{empirical saturation region}, within the family
of prompts we explored, as the accuracy range in which further prompt
iterations produced changes that did not consistently generalize
across problem distributions (Section~\ref{sec:official}), holding
inference parameters fixed. The two prompt variants evaluated at full
scale ($n=400$) scored 71.8\% (AN38; 95\%~CI: [67.1\%, 75.9\%]) and
79.3\% (AN45c; 95\%~CI: [75.0\%, 82.9\%]). We use these as the
full-scale anchors of the region analyzed here. Smaller-sample
evaluations of additional variants shown in Table~\ref{tab:single},
which carry wider confidence intervals, provide broader context but do
not independently establish a saturation boundary.

Among variants for which Table~\ref{tab:single} reports
class-conditional recalls, we observe an apparent trade-off: variants
with higher \textsc{True} recall tend to have lower \textsc{False}
recall, and vice versa. We do not construct a formal Pareto frontier
across variants evaluated at different sample sizes ($n=50$ to
$n=400$), since smaller-$n$ estimates carry substantially wider
confidence intervals that would make any claimed dominance relation
unreliable. At the two variants evaluated at full scale, the trade-off
is visible directly: AN45c achieves higher \textsc{True} recall than
AN38 (95.9\% vs.\ 78.5\%) at a slightly lower \textsc{False} recall
(63.4\% vs.\ 65.4\%).

\subsection{The Merge Pattern: $\mathrm{avg}(A,B)$, Not $\max(A,B)$}
\label{sec:mergepattern}

AN35 and AN35b are complementary specialists: AN35 achieves
\textsc{True}=58.3\%, \textsc{False}=84.6\%; AN35b achieves
\textsc{True}=79.2\%, \textsc{False}=65.4\%. Their arithmetic means
are \textsc{True}=68.8\%, \textsc{False}=75.0\%. The merge result AN38
produces \textsc{True}=70.8\%, \textsc{False}=76.9\%~---~within~2pp
of the arithmetic mean on both dimensions, not near the respective
maxima. Three subsequent attempts (AN45d, AN45e, AN45f) replicated the
averaging result. Combining complementary prompts in a static text file
yields average performance, not maximum performance.

\subsection{Routing Was Not Achieved by Any Prompt in This Family}

Comparing AN45c and AN38~---~the two variants we evaluated at full
scale~---~illustrates a contrast: AN45c, which places the trivial-magma
check before the counterexample table, shows higher \textsc{True}
recall and slightly lower \textsc{False} recall than AN38, which
places the counterexample table first (Table~\ref{tab:single}). We did
not systematically code all 45+ variants by strategy emphasis, so we
report this as a contrast between these two variants rather than a
general pattern across the family. We did not identify, among the
variants we tested, a prompt whose observed predictions demonstrated
reliable per-instance routing~---~applying one strategy for equations
with a specific syntactic form and the other strategy otherwise. When
we constructed a prompt from two complementary parents (AN38, from
AN35 and AN35b), the result's recall profile was close to the
arithmetic mean of its parents rather than close to either parent's
maximum (Section~\ref{sec:mergepattern}); three subsequent merge
attempts showed the same pattern.

We report this as an empirical pattern within our explored family. We
did not measure attention weights or other internal signals, so we do
not claim a general mechanism for how language models execute prompt
instructions.

\subsection{Theoretical Upper Bound and the Path Beyond}

gpt-oss-120b with no cheatsheet achieves~91.7\% \textsc{True} recall
and~92.3\% \textsc{False} recall on \emph{normal} problems. If a
perfect external router could direct each hard problem to the
appropriate strategy, the theoretical ceiling would be approximately:
\[
  0.5 \times 91.7\% + 0.5 \times 92.3\% = 92.0\%
\]
---~well above the observed~71--79\% saturation region. Gemini~1.5~Pro
achieves~90.2\% on hard problems with no cheatsheet, suggesting that
frontier models have internalized routing-equivalent algebraic reasoning
that cannot be injected through prompting into weaker models.

Concretely, escaping the empirical saturation region for gpt-oss-120b
on this task likely requires one of:
(a)~an ensemble of specialized prompts with external routing by problem
type; (b)~fine-tuning on the Equational Theories Project graph
($\approx$22~million labeled equation pairs); or~(c)~a hybrid
LLM~+~symbolic verifier architecture where the LLM proposes candidate
counterexamples and a Mace4-equivalent tool verifies them.

% ─────────────────────────────────────────────
\section{Official Benchmark Validation}
\label{sec:official}

The complete competition leaderboard provides an independent validation
dataset. We use it only to evaluate whether the empirical patterns
identified in our controlled ablation study
(Sections~\ref{sec:results}--\ref{sec:multimodel}) generalize beyond our
own prompts. The analytical sequence is deliberately unidirectional: our
experiments produced the findings; the findings motivated the hypotheses
of Section~\ref{sec:analysis}; the leaderboard, released after our
submission was locked, then serves as an out-of-sample test of whether
those hypotheses are compatible with the behavior of 309 independently
developed submissions.

\subsection{Data and Scope}
\label{sec:official-data}

The final competition leaderboard~\citep{sair2026leaderboard} reports
results for 310~evaluated submissions of 1{,}007 registered participants,
across two evaluation tracks. The \emph{Evaluation (Scored)} board, which
determined official rankings, averages accuracy over three models
(GPT-OSS~120B, Gemma~4~31B~IT, Llama~3.3~70B~Instruct) and three problem
sets (\texttt{normal}, \texttt{hard}, \texttt{extra\_hard}, 200~problems
each, three runs per problem). The \emph{Order-5 (Research)}
board~\citep{sair2026leaderboardresearch} evaluates the same three models
on a separate 200-problem set restricted to order-5
magmas.\footnote{Because the two boards draw from different problem
distributions, the pair (scored accuracy, research accuracy) provides,
for every submission, a cross-distribution measurement under a fixed
prompt.}

Because both boards were published after the submission deadline, no
participant---ourselves included---could optimize against them. The
leaderboard therefore functions as a natural pre-registered test: every
hypothesis stated in Sections~\ref{sec:analysis}
and~\ref{sec:multimodel} of the versions of this paper published before
competition close (v13--v15.2, timestamped on arXiv and Zenodo) predates
the evidence examined here.

Summary statistics frame the difficulty of the task. On the scored
board, mean accuracy across all 310~submissions is 55.3\% (median
55.8\%), and the maximum is 69.1\%. On the research board, the mean is
61.3\% and the maximum 82.4\%.

An earlier Contributor Network snapshot ($n{=}52$; April~20, 2026)
reported dual hard2/hard3 scores for 11 external submissions
(Table~\ref{tab:crossdist}). That snapshot already exhibited large
cross-split collapses---Betka's 99.0\% on hard2 falls to 56.3\% on
hard3 ($-$42.7pp); Woon Siang Yi's 81.3\% on hard3 falls to 51.0\% on
hard2 ($-$30.3pp), with the participant labeling the submission
\texttt{hard3\_overfitted}---and flagged Arjun Garg's
\texttt{bank\_lookup\_v8} (83.0\% hard2, 70.8\% hard3) as the sole
submission above 65\% on both splits. We treat this as a historical
precursor: the full-leaderboard test in
Section~\ref{sec:official-tradeoff} supersedes it at scale.

\begin{table}[h]
\centering
\caption{Cross-distribution performance: hard2 vs.\ hard3
(GPT-OSS 120B, SAIR official benchmark). Contributor Network
external submissions with dual-split scores as of
April~20, 2026~\citep{sair2026leaderboard} ($n{=}11$; own
submissions omitted; see Section~\ref{sec:local-official}).
Participants cited by name under SAIR's open science framework.}
\label{tab:crossdist}
\small
\begin{tabular}{llrrrr}
\toprule
Participant & Cheatsheet & hard2 Acc. & hard3 Acc. & $|\Delta|$ & hard3 F1 \\
\midrule
Betka       & 98\_hard200           & \textbf{99.0\%} & 56.3\% & 42.7pp & 55.9\% \\
Pandey      & hard98\_derivate\_2   & 97.0\%          & 55.5\% & 41.5pp & 53.9\% \\
Reza Jamei  & traditional\_ml       & 89.0\%          & 61.3\% & 27.7pp & 68.6\% \\
Reza Jamei  & combo\_jack           & 86.0\%          & 58.5\% & 27.5pp & 66.5\% \\
Arjun Garg  & bank\_lookup\_v8      & \textbf{83.0\%} & \textbf{70.8\%} & 12.2pp & \textbf{74.6\%} \\
Heath       & distilled-r-12        & 80.0\%          & 59.0\% & 21.0pp & 62.6\% \\
Woon Siang Yi\textsuperscript{$\dagger$} & hard3\_overfitted & 51.0\% & \textbf{81.3\%} & 30.3pp & \textbf{83.1\%} \\
Garg        & bank\_lookup\_v5      & 51.5\%          & 60.8\% & 9.3pp  & 69.3\% \\
SimonRJ     & Stage 1 Prompt        & 58.0\%          & 58.3\% & 0.3pp  & 67.3\% \\
Debtirtha Saha & Test 8             & 60.0\%          & 71.0\% & 11.0pp & 76.0\% \\
Devanshu Dixit & chat\_v32          & 49.5\%          & 68.0\% & 18.5pp & 74.3\% \\
\bottomrule
\multicolumn{6}{l}{\textsuperscript{$\dagger$}Self-labeled ``hard3\_overfitted'' by participant.}
\end{tabular}
\end{table}

\subsection{Cross-Distribution Trade-Off at Full Scale}
\label{sec:official-tradeoff}

Version~15.2 of this paper, using the voluntary Contributor Network
snapshot ($n{=}52$, 11 external submissions with dual-split scores),
observed that only one submission exceeded 65\% accuracy on both hard2
and hard3, and flagged the full leaderboard as the natural
falsification test of whether that dual-split robustness would persist
at scale. The complete leaderboard resolves this test on the
scored--research axis analyzed below.

\begin{table}[t]
\centering
\caption{Cross-board performance thresholds, final leaderboard
($n{=}310$). Each submission is evaluated on both the official scored
distribution and the order-5 research distribution.}
\label{tab:dualthreshold}
\begin{tabular}{lc}
\toprule
Criterion & Submissions \\
\midrule
$>$60\% on both boards & 27 / 310 \\
$>$65\% on both boards & 5 / 310 \\
$>$70\% on both boards & \textbf{0 / 310} \\
\bottomrule
\end{tabular}
\end{table}

No submission among all 310 exceeds 70\% accuracy on both distributions
(Table~\ref{tab:dualthreshold}). The pattern is most visible at the
top of the research board: the five highest-scoring research submissions
lose between 19 and 32~percentage points when evaluated on the official
scored distribution (Table~\ref{tab:researchcollapse}).

\begin{table}[t]
\centering
\caption{The five highest-ranked submissions on the research board and
their official scored results. Team names are public leaderboard
entries~\citep{sair2026leaderboard, sair2026leaderboardresearch}.}
\label{tab:researchcollapse}
\begin{tabular}{lccc}
\toprule
Team & Research & Scored & Scored rank \\
\midrule
Weihang Ding & 82.4\% & 56.3\% & 137 \\
Lothar Schiemanowski & 81.9\% & 50.3\% & 288 \\
Math Distillation & 80.6\% & 61.6\% & 20 \\
Ken & 80.4\% & 56.2\% & 141 \\
zkdmt & 80.1\% & 52.2\% & 264 \\
\bottomrule
\end{tabular}
\end{table}

The exception identified in v15.2 also resolves. Arjun Garg's
submission, which exceeded 65\% on both hard2 and hard3 in the
Contributor Network snapshot and which we flagged as a potential
falsifier, ranks 8th on the research board
(77.9\%) but 166th on the scored board (55.5\%)---below our own
submission's scored accuracy. What appeared at $n{=}11$ to be an escape
from the dual-split pattern is, at $n{=}310$, another instance of
cross-board fragility.

Conversely, the five submissions that do clear 65\% on both boards
(Table~\ref{tab:robustfive}) are characterized not by peak accuracy on
either distribution but by small cross-board gaps. The competition
winner exhibits the smallest gap of the top group
($+1.0$pp), despite ranking only 15th--20th on the research board. At
full scale, cross-distribution robustness---not peak accuracy on any
single distribution---predicts official ranking.

\begin{table}[t]
\centering
\caption{All submissions exceeding 65\% on both boards ($n{=}5$ of
310), ordered by scored accuracy.}
\label{tab:robustfive}
\begin{tabular}{lccc}
\toprule
Team & Scored & Research & Gap \\
\midrule
Dufius & 69.1\% & 70.1\% & $+1.0$ \\
vt & 67.9\% & 66.2\% & $-1.7$ \\
Reza Jamei & 66.3\% & 74.1\% & $+7.8$ \\
AJ & 66.2\% & 70.2\% & $+4.0$ \\
yan-biao & 65.3\% & 74.3\% & $+9.0$ \\
\bottomrule
\end{tabular}
\end{table}

Two conclusions follow for this external evidence. First, under a
homogeneous-linear null calibrated to the full sample
(Pearson~$r{=}0.74$; 10{,}000 bootstrap replicates, seed~42),
research-selected upper subsets fall outside the expected correlation
interval at all three cuts examined (top~33\%/~25\%/~10\%), whereas
scored-selected upper subsets do not do so consistently---an observed
directional association on the research$\to$scored axis, not a claim of
shared mechanism with the local--official divergence of
Section~\ref{sec:local-official}~\citep{cazares2026sair}. Second, the
empirical saturation region identified in Section~\ref{sec:ceiling} for
our prompt family (approximately 71.8--79.3\% on local hard3 at full
scale) is consistent with the official ceiling observed across
participants: no submission exceeds 69.1\% on the official scored
evaluation, and no submission exceeds 70\% on both boards
simultaneously.

\subsection{Implications for the Saturation Region}

These results reframe the saturation region finding. The ceiling is not
merely an upper bound on accuracy~---~it is a \emph{fragility zone}
where gains are distribution-dependent and may fail to transfer across
evaluation distributions.

For practitioners, this implies: a cheatsheet that substantially
outperforms baseline on a held-out set does not guarantee
generalization. Cross-distribution validation~---~running the same
variant on two structurally different problem samples~---~is necessary
to distinguish genuine ceiling improvements from distribution-specific
optimization.

\subsection{Compatibility with the Structural Classification Hypothesis}
\label{sec:official-router}

Section~\ref{sec:analysis} advanced the hypothesis---prior to the
release of any official results---that gpt-oss-120b behaves on this
task as a structural-pattern classifier rather than a prover, and that
prompts succeed to the extent that they support structural
classification rather than open-ended mathematical reasoning. The
leaderboard offers an observational test: if the hypothesis is
directionally correct, strategies that explicitly constrain the model's
decision process should be over-represented at the top of the official
ranking.

The evidence is compatible with the hypothesis. The winning
submission's publicly listed category places it as a rule-based
classifier rather than a reasoning-guidance prompt, and among the
five submissions exceeding 65\% on both boards it shows the smallest
cross-board gap (Table~\ref{tab:robustfive}).\footnote{We cite only
the publicly visible category and published scores of leaderboard
submissions. A content-level analysis of other participants' prompts
is outside the scope of this paper.} We treat this as external
supporting evidence, not as proof: a single winning instance cannot
establish that decision-space reduction is necessary for escaping the
saturation region, and the leaderboard does not permit controlled
comparison of strategy categories. One publicly documented winning
strategy is consistent with the prediction, formulated in
Section~\ref{sec:analysis} before official results existed, that
strategies constraining the model's decision process would outperform
reasoning guidance.

A second, categorically distinct submission is consistent with the
same interpretation: one participant (Heath, \texttt{distilled-rules-12})
is publicly categorized as a structural classifier rather than a
CE-table or reasoning-guidance approach, achieving~80.0\% on hard2
and~59.0\% on hard3---competitive with CE-table approaches in our
comparison set. We note this as a second independent observation
consistent with the hypothesis that structural classification
strategies are competitive on this task; we do not characterize its
internal design further.

Notably, one participant (McKenna, unpublished) reported achieving
approximately~70\% accuracy across all three evaluation models
simultaneously using a single theory-grounded cheatsheet~---~the only
submission observed to maintain consistent performance across
gpt-oss-120b, Llama~3.3~70B, and Gemma~4~31B. This convergence across
models with substantially different capacity profiles suggests that
mathematically grounded approaches may generalize more robustly than
empirically iterated CE-table approaches, a hypothesis warranting
systematic investigation in future work.


Section~\ref{sec:instability} develops the observable quantity that,
in our view, underlies this pattern: the degree to which a prompt
delegates the classification decision to the model, measurable through
run-to-run verdict instability.

% ─────────────────────────────────────────────
\section{Run Instability and Decision Delegation}
\label{sec:instability}

The official evaluation executes every problem three times per model.
The public leaderboard exposes these per-run outcomes at the problem
level, making it possible to measure a quantity that aggregate accuracy
conceals: how often a model returns \emph{different verdicts across
identical runs} of the same problem under the same prompt. We call the
fraction of problems receiving inconsistent verdicts the
\textbf{disagreement rate}, and we propose it as an observable proxy
for \emph{decision delegation}: the degree to which a prompt leaves the
classification decision to the model's own stochastic reasoning process
rather than constraining it procedurally.

\subsection{Data and Validation}
\label{sec:instability-data}

We extracted per-problem, per-run outcomes from the public leaderboard
problem matrices for two submissions spanning the strategy spectrum:
our own (AN45c, 2{,}252~bytes) and the winning submission
(9{,}834~bytes; Section~\ref{sec:official-router}). For each submission
this yields $200 \times 3 = 600$ run-level outcomes per
model$\times$set cell---5{,}400 outcomes across the scored board's
nine cells per submission, plus 1{,}800 for our submission on the
research board.

To validate the extraction, we reconstructed each cell's official
accuracy from the run-level data. Twenty of twenty-one model$\times$set
combinations reproduce the officially published accuracy exactly; the
remaining combination is a partial extraction and matches on its
subset.\footnote{Extraction scripts and the validated dataset are
released in the companion repository~\citep{cazares2026sair}.}

\subsection{Disagreement Is Prompt-Dependent and Inverts Across Models}
\label{sec:instability-results}

Table~\ref{tab:disagreement} reports disagreement rates for both
submissions across all models and problem sets.

\begin{table}[t]
\centering
\caption{Disagreement rate: fraction of problems (of 200 per cell)
receiving inconsistent verdicts across three identical runs
(temperature~0). ``Ours'' is AN45c; ``Winner'' is the first-place
submission. GPT-OSS normal for our submission is measured on the
extracted subset ($n{=}200$, combining both extraction passes).}
\label{tab:disagreement}
\small
\begin{tabular}{llrr}
\toprule
Model & Problem set & Ours & Winner \\
\midrule
GPT-OSS 120B & normal      & 39.0\% & 8.5\% \\
GPT-OSS 120B & hard        & \textbf{50.5\%} & 9.0\% \\
GPT-OSS 120B & extra\_hard & 43.0\% & 5.0\% \\
\midrule
Gemma 4 31B  & normal      & 16.0\% & 12.0\% \\
Gemma 4 31B  & hard        & 10.5\% & 22.5\% \\
Gemma 4 31B  & extra\_hard & 28.0\% & 10.0\% \\
\midrule
Llama 3.3 70B & normal      & 2.0\% & 48.0\% \\
Llama 3.3 70B & hard        & 1.0\% & \textbf{54.5\%} \\
Llama 3.3 70B & extra\_hard & 1.0\% & 45.5\% \\
\bottomrule
\end{tabular}
\end{table}

Two patterns stand out. First, on GPT-OSS~120B---the model where both
prompts do their primary work---the winning submission's disagreement
rate is five to nine times lower than ours. Under our heuristic
guidance prompt, half of all hard problems (50.5\%) receive
inconsistent verdicts across three identical runs; under the winning
submission, only 9\% do. The winning submission exhibits
substantially lower run-to-run disagreement than our
reasoning-guidance prompt on this model. This pattern is consistent
with the interpretation that the winning submission constrains the
model's decision process more tightly than reasoning guidance does:
under the winning submission, identical runs converge on identical
verdicts far more often; under ours, the model more often resolves
the same problem differently run to run.

Second, the pattern \emph{inverts} on Llama~3.3~70B. Our prompt makes
Llama almost perfectly deterministic (1--2\% disagreement)---but
deterministic in a degenerate way. The per-run data reveals the exact
mechanism of the Llama collapse documented in
Section~\ref{sec:multimodel}: on every problem set, our submission
produces approximately 98--100 solid-correct and 98--100 solid-incorrect
cells with near-zero \textsc{True}-class F1. On a balanced set, this is
the signature of a constant classifier: Llama answers \textsc{False}
on essentially every problem, three runs out of three---correct on the
\textsc{False} half, wrong on the \textsc{True} half. The winner's
prompt, by contrast, destabilizes Llama (45--55\% disagreement) but
makes it \emph{useful}: 63.0\% accuracy on hard versus our 50.5\%.
Decision delegation is therefore not uniformly good or bad; its
observed effect depends on what the model does with the delegated
decision. GPT-OSS's disagreement is concentrated on unresolved
problems, while Llama's near-total agreement reflects convergence on a
single, largely incorrect class.

\subsection{Instability Is Not Noise That Voting Can Remove}
\label{sec:instability-voting}

If run-to-run disagreement were symmetric noise around a correct
signal, majority voting over the three runs would recover accuracy.
It does not. For our submission on GPT-OSS hard, run-level accuracy is
53.8\% and majority-vote accuracy is 53.0\%---no improvement. The
reason is visible in the composition of the disagreeing cells: of the
101 mixed cells, 53 have one correct run and 48 have two, an almost
uniform split. Under our prompt, the model is not noisily orbiting the
right answer on these problems; it is at chance on them. The
disagreement rate thus identifies the subpopulation of problems on
which the prompt fails to constrain the model to \emph{any} decision,
which is qualitatively different from failing to constrain it to the
\emph{correct} one.

\subsection{Content, Not Length, Drives the Llama Collapse}
\label{sec:instability-llama}

Earlier versions of this work attributed the Llama collapse to prompt
length, based on the within-family observation that our prompts
exceeding $\approx$2KB produced near-zero \textsc{True} recall. The
full leaderboard refutes length as the driver. Across all
310~submissions, 179 exhibit the Llama collapse signature
(\textsc{True}-class F1 below 5\% on the normal set), and their mean
cheatsheet size (4.2KB) is \emph{smaller} than that of the
131~submissions where Llama remains functional (6.1KB). Our own
2{,}252-byte submission collapses Llama; the winner's 9{,}834-byte
submission does not. Within our design family, length covaried with
rule-system complexity, producing a spurious length effect; at
leaderboard scale, the collapse is evidently a function of prompt
\emph{content}---which instructions the prompt gives and how they
interact with Llama's instruction-following capacity---rather than of
size. We accordingly retract the size-based formulation and restate
the finding as content-dependent: certain prompt structures, including
our CE-table family, drive Llama~3.3~70B into a deterministic
single-class policy regardless of their length.

\subsection{Interpretation and Scope}
\label{sec:instability-interp}

These observations connect the empirical saturation region
(Section~\ref{sec:ceiling}) to the structural classification
hypothesis (Sections~\ref{sec:analysis}
and~\ref{sec:official-router}) through a measurable intermediate
variable. Under prompts in our heuristic-guidance family, run-to-run
inconsistency is high and concentrated exactly on the problems the
model cannot reliably solve---consistent with the interpretation that
these prompts leave a large fraction of the decision to the model's
stochastic reasoning. The winning submission is consistent with the
same interpretation from the opposite direction: on GPT-OSS~120B it
exhibits disagreement rates five to nine times lower than our
heuristic-guidance prompt. We conjecture---but with $n{=}2$ prompts
cannot establish---that disagreement rate is a general leading
indicator of distribution fragility: a prompt whose accuracy rests on
a large population of at-chance problems has more room to lose under
distribution shift than one whose accuracy rests on deterministically
solved problems. Testing this conjecture requires per-run data for a
larger sample of submissions, which the public leaderboard makes
possible; we leave this to future work.

Two scope limitations apply. The contrast in
Table~\ref{tab:disagreement} spans two prompts, chosen to represent
the endpoints of the observed strategy spectrum; intermediate designs
are unmeasured. And the disagreement rate is measured under the
official pipeline's sampling configuration; its magnitude (though
plausibly not its ordering) may vary under other configurations.

% ─────────────────────────────────────────────
\section{Conclusion}
\label{sec:conclusion}

\paragraph{Summary of contributions.}
We present the first systematic empirical characterization of the
single-prompt ceiling in formal reasoning tasks. Across~45+ prompt
variants, balanced hard accuracy for gpt-oss-120b on hard3 plateaus at
$\approx$71--79\% at scale. Merge experiments confirm that combining
complementary prompts produces accuracy near the arithmetic mean of the
parents, not near their maxima~---~a structural constraint, not a local
optimum. Our highest-performing local submission (AN45c, 2{,}252~bytes)
achieves~79.25\% on gpt-oss-120b ($n=400$; 95\%~CI: [75.0\%, 82.9\%]),
with \textsc{True} recall of~95.9\% and \textsc{False} recall of~63.4\%.
AN38 is our most robust submission under distribution shift, producing
$+$5.6pp over the official baseline (Section~\ref{sec:official}). A
cross-provider validation run on OpenRouter/DeepInfra~bf16 ($n=20$)
yielded~90--95\%, consistent with the full-scale local result. The token
cap finding demonstrates that Gemma's performance (50\% vs.~85\%) is
entirely determined by provider configuration, not model capability.
Model-specific bias profiles~---~\textsc{True} bias in gpt-oss-120b,
instruction-capacity collapse in Llama, token-gated reasoning in
Gemma~---~are stable across variants and not correctable by prompt
engineering alone.

Post-submission analysis reveals that AN45c (79.25\% local,
$n=400$) scores~55.5\% on the official hard3 benchmark~---~4.3pp below
the no-cheatsheet baseline~---~while AN38 produces a robust $+$5.6pp
improvement (65.3\%) with the highest balanced F1 score (67.6\%) among
non-overfit submissions in the Contributor Network ($n=52$ visible
submissions, April~20, 2026; Section~\ref{sec:local-official}). On the
full leaderboard (Section~\ref{sec:official-tradeoff}), no submission
exceeds 70\% on both the scored and research boards, and
research-selected upper subsets fall outside the homogeneous-null
correlation interval at all three cuts examined.

\paragraph{Limitations.}
The primary AN45c result (79.25\%, $n=400$) is statistically reliable
(95\%~CI: [75.0\%, 82.9\%]); however, its \textsc{True} recall~(95.9\%)
and \textsc{False} recall~(63.4\%) reflect an imbalanced profile that
may not generalize across providers or problem distributions. Small-sample
limitations remain for Gemma (85\%, $n=20$) and DeepSeek (80\%, $n=10$),
both carrying $\pm$10--15pp confidence intervals. The cross-provider
OpenRouter/DeepInfra run ($n=20$, 90--95\%) should be read as a
consistency check, not an independent measurement. The Gemma official
provider (OpenRouter/Novita~bf16) suppresses reasoning mode; we were
unable to test this configuration in a controlled way. All results are
specific to equational implication over magmas; generalization to other
formal reasoning domains is an open question.

The primary AN45c result (79.25\%, $n=400$, Together~AI~bf16) does not
generalize to the SAIR official benchmark (55.5\%, $n=20$,
DeepInfra~bf16), a gap of~$-$23.75pp. AN38's smaller gap ($-$6.5pp
local to official) suggests that simpler prompts generalize more
robustly within the saturation region. An earlier Contributor Network
snapshot ($n{=}52$ of 1{,}007 participants) was a voluntary public
sample; Section~\ref{sec:official} analyzes the final competition
leaderboard ($n{=}310$).

\paragraph{Future work.}
Three directions follow from the ceiling characterization: (1)~an
ensemble of two specialized prompts with external routing by syntactic
problem features; (2)~fine-tuning on the Equational Theories Project
implication graph~($\approx$22M labeled pairs); and~(3)~replication
on other computationally asymmetric formal reasoning tasks
(satisfiability checking, reachability in formal systems).

\paragraph{Closing.}
The central lesson is structural: what improved performance was not
teaching the model new mathematics but controlling the order in which
the model applies the mathematics it already knows. Expanding the
model's reasoning repertoire through longer, more elaborate prompts
consistently underperformed constraining its reasoning flow through
minimal, well-ordered instructions. In formal reasoning as in
engineering: less, structured deliberately, beats more.

% ─────────────────────────────────────────────
\bibliographystyle{plainnat}
\bibliography{references}

\end{document}
